\documentclass[conference]{IEEEtran}

\usepackage{graphicx}
\usepackage{amsmath}
\usepackage{float}

\title{Design and Implementation of a Line Following Robot Using Discrete Transistor Logic and TIP122 Darlington Transistors}



\begin{document}

\author{
\centering
\begin{minipage}[t]{0.32\linewidth}
\centering
Umer Faisal \\
\textit{Faculty of Computer Science and Engineering \\
Ghulam Ishaq Khan Institute of Engineering Sciences and Technology \\
Topi, Pakistan} \\
u2024647@giki.edu.pk
\end{minipage}%
}

\maketitle

\begin{abstract}
This paper presents the design and implementation of a line-following robot using infrared reflectance sensors and discrete electronic components without the use of microcontrollers or programmable logic devices. The system employs TCRT5000 infrared sensors for line detection and transistor-based signal processing for decision making. Motor control is achieved using TIP122 Darlington transistors operating as low-side switches to drive DC motors.

The design focuses on understanding fundamental electronic principles such as transistor switching, signal inversion, inductive load behavior, and power isolation. Key challenges encountered during implementation include active-low sensor interpretation, loading effects between transistor stages, grounding issues between multiple power supplies, and inductive voltage spikes generated by DC motors. These issues were resolved through careful circuit design, proper biasing, and the use of flyback diode protection.

Experimental results demonstrate that the robot is capable of successfully following a predefined black line on a contrasting surface by dynamically controlling motor states based on sensor input. The project highlights the effectiveness of discrete component-based design for educational robotics and provides insight into real-world electronic system integration.
\end{abstract}

\begin{IEEEkeywords}
Line Following Robot, TIP122, Transistor Logic, TCRT5000, Embedded Systems, Analog Electronics, Robotics
\end{IEEEkeywords}

\section{Introduction}

Autonomous mobile robots are widely used in industrial automation, navigation systems, and educational platforms. Among these, line-following robots represent one of the simplest yet most effective systems for demonstrating the principles of feedback control, sensor integration, and motor actuation.

In modern implementations, microcontrollers such as Arduino or PIC are commonly used to process sensor data and control motors. However, this project explores a hardware-centric approach in which all logic and control operations are implemented using discrete electronic components. This approach provides deeper insight into the underlying principles of electronic circuit behavior, including transistor switching characteristics, signal conditioning, and power electronics.

The primary objective of this project is to design and implement a fully functional line-following robot capable of detecting and tracking a black line on a contrasting surface using infrared sensors. The system uses TCRT5000 sensors to detect surface reflectance variations and converts these signals into digital logic levels using transistor-based inverter circuits. The processed signals are then used to control DC motors via TIP122 Darlington transistors configured as switching devices.

A key challenge in such systems is the correct interpretation of active-low sensor outputs, which requires signal inversion before motor actuation. Additionally, since motors are inductive loads, they introduce back electromotive force (back EMF) during switching events, which can damage electronic components if not properly managed. These challenges are addressed in this work through appropriate circuit design techniques, including the use of pull-up resistors, common grounding between power supplies, and flyback diode protection across motor terminals.

\section{Project Requirements}

The development of the line-following robot required the integration of sensing, signal processing, motor driving, and mechanical motion subsystems. The primary requirement of the project was to design a robot capable of autonomously detecting and following a predefined black line on a contrasting white surface without the use of a microcontroller.

The system required infrared sensors for surface detection, transistor-based logic circuitry for signal inversion and processing, and power transistor drivers for motor actuation. Separate power supplies were required for the logic circuitry and DC motors to ensure stable operation under varying load conditions. Additionally, protection circuitry was necessary to safeguard the motor drivers from inductive voltage spikes generated by the motors.

From a mechanical perspective, the robot required a lightweight chassis, sufficient wheel traction, and stable differential drive movement to achieve reliable line tracking and turning performance. The complete system was designed using low-cost and easily available electronic components to ensure simplicity, affordability, and ease of implementation.

\section{System Architecture}


The overall architecture of the proposed line-following robot is shown in Fig.~\ref{fig:block}. The system is composed of multiple interconnected subsystems responsible for sensing, signal processing, motor driving, and mechanical motion. The robot operates using a differential drive mechanism in which two independently controlled DC motors determine the direction and movement of the robot.

The primary sensing elements used in the system are TCRT5000 infrared reflectance sensors. These sensors continuously monitor the surface beneath the robot and detect the presence or absence of the black tracking line based on infrared light reflection. White surfaces reflect infrared radiation more effectively, whereas black surfaces absorb most of the emitted light. As a result, the sensor output changes depending on the detected surface.

The outputs of the sensors are active-low, meaning that the sensor produces a LOW voltage when the white surface is detected and a HIGH voltage when the black line is present. Since the motor driver stage requires active-high control signals for proper operation, transistor-based inverter circuits are introduced between the sensors and the motor drivers. These inverter stages convert the active-low sensor outputs into active-high which is suitable for motor actuation.

The signal conditioning stage is implemented using 2N3904 NPN transistors configured in common-emitter mode. Pull-up resistors are connected to the collector terminals to ensure valid HIGH output levels when the transistor is in the OFF state. This stage also isolates the sensor subsystem from the motor driver stage and prevents loading effects that could distort signal levels.

Motor control is achieved using TIP122 Darlington transistors operating as low-side switches. Each transistor controls one DC motor independently. When the base of the TIP122 receives a HIGH signal from the inverter stage, the transistor enters saturation and provides a conductive path between the motor and ground, allowing current to flow through the motor. When the base signal is LOW, the transistor switches OFF and the motor stops.

To protect the transistor drivers from inductive voltage spikes generated by the DC motors, flyback diodes are connected across each motor terminal. These diodes provide a safe path for inductive current during switching events and prevent damage caused by back electromotive force (back EMF).

The robot uses two separate power supplies: a 5V supply for the sensors and transistor logic circuitry, and a 9V battery for the DC motors. Although separate voltage domains are used, both systems share a common ground reference. This common grounding is essential for correct transistor switching and stable operation of the complete system.

The mechanical subsystem consists of a four-wheel chassis with two front motor-driven wheels and two rear support wheels. Steering is achieved through differential motion, where the relative speeds of the left and right motors determine the turning direction of the robot. When both motors operate simultaneously, the robot moves forward, whereas selective motor activation enables directional correction during line tracking.

The modular architecture of the system simplifies debugging, testing, and future expansion. Each subsystem can be analyzed independently, allowing systematic identification and correction of electrical and mechanical issues encountered during implementation.

\section{Proteus Simulation}


Before hardware implementation, the complete circuit was designed and tested in \textit{Proteus Design Suite}. The simulation environment was used to verify the operation of the sensor logic, transistor inverter stages, TIP122 motor drivers, and flyback protection circuitry. 

\section{Hardware Components Specification}

\begin{table}[h]
\centering
\caption{Hardware Components and Specifications}
\label{tab:hardware}

\resizebox{\columnwidth}{!}{%

\begin{tabular}{|c|l|l|l|}
\hline
\textbf{Sr. No.} & \textbf{Component} & \textbf{Specification} & \textbf{Purpose} \\ \hline

1 & TCRT5000 IR Sensor Module &
5V Active-Low Reflective Sensor &
Line detection \\ \hline

2 & 2N3904 Transistor &
NPN Bipolar Junction Transistor &
Signal inversion and logic conditioning \\ \hline

3 & TIP122 Transistor &
NPN Darlington Power Transistor &
DC motor driving \\ \hline

4 & DC Motor &
9V Geared DC Motor &
Robot movement \\ \hline

5 & Flyback Diode &
1N4007 Rectifier Diode &
Back EMF protection \\ \hline

6 & Resistors &
1k$\Omega$, 10k$\Omega$ &
Current limiting and biasing \\ \hline

7 & Power Supply (Logic) &
5V DC Supply &
Sensor and logic power \\ \hline

8 & Power Supply (Motor) &
9V Battery &
Motor power source \\ \hline

9 & Chassis &
Four-wheel Acrylic Chassis &
Mechanical structure \\ \hline

10 & Wheels &
Rubberized Plastic Wheels &
Traction and motion \\ \hline

11 & Breadboard &
Standard Solderless Breadboard &
Circuit prototyping \\ \hline

12 & Connecting Wires &
Jumper Wires &
Electrical connections \\ \hline

\end{tabular}

}
\end{table}

\section{System Design}

The proposed line-following robot is composed of multiple interconnected subsystems responsible for sensing, signal processing, motor driving, power distribution, and mechanical locomotion. Each subsystem performs a specific task that contributes to the overall operation of the robot. The modular architecture simplifies implementation, debugging, and future upgrades.


\subsection{Sensor Subsystem}

The sensor subsystem is responsible for detecting the black tracking line on the surface and generating corresponding electrical signals for motor control. The robot uses two TCRT5000 infrared reflectance sensor modules mounted near the bottom-front side of the chassis. These sensors continuously monitor the surface beneath the robot and determine whether the sensor is positioned above the black line or the surrounding white surface.

The TCRT5000 sensor module operates using the principle of infrared reflection. The module contains an infrared emitting diode (IR LED) and a phototransistor receiver positioned side by side inside the sensor package. The infrared LED continuously emits infrared radiation toward the ground surface beneath the robot. Depending on the reflectivity of the surface, different amounts of infrared light are reflected back toward the phototransistor.

White surfaces reflect a large amount of infrared radiation. When the sensor is placed above a white surface, most of the emitted infrared light is reflected back to the phototransistor receiver. This causes the phototransistor to conduct more current. Conversely, black surfaces absorb most of the incident infrared radiation and reflect very little light back to the sensor. As a result, the phototransistor conduction decreases significantly when the sensor is positioned above the black line.

The phototransistor generates an analog electrical response proportional to the intensity of reflected infrared light. However, analog signals are highly sensitive to environmental noise, ambient lighting conditions, and surface variations. To produce stable digital outputs suitable for motor control logic, the sensor module includes an onboard voltage comparator integrated circuit, typically the LM393 comparator.

The LM393 comparator compares the analog voltage generated by the phototransistor with a user-adjustable reference voltage. One input terminal of the comparator receives the sensor signal produced by the phototransistor, while the other input terminal receives a reference threshold voltage generated using an onboard potentiometer.

The comparator continuously evaluates the relationship between these two voltages:

\begin{itemize}
    \item If the sensor voltage exceeds the reference voltage, the comparator changes the output state.
    \item If the sensor voltage falls below the reference voltage, the output switches to the opposite logic state.
\end{itemize}

The potentiometer allows adjustment of the reference threshold voltage and therefore controls the sensitivity of line detection. This enables calibration for different lighting conditions, surface materials, sensor heights, and line contrast levels.

The sensor module used in this project exhibits active-low behavior. Therefore, when the sensor detects a white reflective surface, the strong infrared reflection causes the comparator output to become LOW. When the sensor detects the black line, weak infrared reflection causes the comparator output to transition HIGH.

The digital logic behavior can therefore be summarized as follows:

\begin{itemize}
    \item White surface detected $\rightarrow$ high infrared reflection $\rightarrow$ LOW output
    \item Black line detected $\rightarrow$ low infrared reflection $\rightarrow$ HIGH output
\end{itemize}

The sensor module also contains current-limiting resistors to protect the infrared LED from excessive current. In addition, onboard indicator LEDs are provided for debugging and calibration purposes. One LED indicates power availability, while another LED reflects the digital output state of the comparator.

The module provides both analog and digital outputs. The analog output directly represents the continuously varying phototransistor signal, whereas the digital output is generated by the LM393 comparator according to the selected threshold voltage. In this project, only the digital output was utilized because the robot required binary logic signals for transistor-based motor control.

Proper sensor placement and calibration are critical for reliable operation. The sensors were mounted close to the ground surface to maximize reflection sensitivity and improve detection accuracy. Experimental testing showed that excessive sensor height reduced reflected infrared intensity and caused unstable detection behavior.

The outputs generated by the sensor modules are forwarded to the signal conditioning and logic inversion stage, where the active-low logic levels are converted into the appropriate signals required for motor actuation.

\subsection{Signal Conditioning and Logic Inversion Subsystem}

Since the TCRT5000 sensors produce active-low outputs, their logic behavior is opposite to the required motor control logic. In the desired operation of the robot, the motor should remain ON when the sensor detects the white surface and should turn OFF when the black line is detected. However, the sensor outputs LOW for white surfaces and HIGH for black surfaces.

To resolve this mismatch, a transistor-based inverter circuit is used. The inverter stage is implemented using an NPN transistor configured in common-emitter mode. A pull-up resistor is connected between the collector terminal and the supply voltage to establish a valid HIGH output level when the transistor is turned OFF.

The operation of the inverter is summarized as follows:

\begin{itemize}
    \item Sensor output LOW (white surface) $\rightarrow$ transistor OFF $\rightarrow$ collector output HIGH
    \item Sensor output HIGH (black line) $\rightarrow$ transistor ON $\rightarrow$ collector output LOW
\end{itemize}

This inversion stage ensures compatibility between the sensor outputs and the motor driver subsystem.

During implementation, loading effects were observed when the motor driver stage was directly connected to the inverter output. The base current drawn by the motor driver transistor reduced the collector voltage of the inverter transistor, preventing a valid HIGH logic level. This issue was resolved by increasing the pull-up resistance and selecting appropriate base resistor values.

\subsection{Motor Driver Subsystem}

The motor driver subsystem controls the operation of the DC motors using TIP122 Darlington transistors. Since the motors require significantly more current than the sensor or logic circuitry can provide directly, transistor-based switching is necessary.

The TIP122 transistor is configured as a low-side switch. The emitter terminal is connected to ground, while the collector terminal is connected to one terminal of the motor. The other motor terminal is connected to the positive terminal of the motor power supply.

When a HIGH signal is applied to the base of the TIP122 transistor through a resistor, the transistor enters saturation mode and allows current to flow through the motor. This causes the motor to rotate. When the base signal becomes LOW, the transistor switches OFF and interrupts current flow through the motor.

The use of a Darlington transistor provides high current gain, allowing small logic-level currents to control comparatively larger motor currents.

\subsection{Protection Subsystem}

DC motors are inductive loads and therefore generate voltage spikes when switched OFF. These voltage spikes, commonly known as back electromotive force (back EMF), can damage transistor junctions and reduce system reliability.

To protect the TIP122 motor driver transistors, flyback diodes are connected across the motor terminals in reverse bias configuration. Under normal operation, the diode remains reverse biased and does not conduct current. However, when the motor is switched OFF, the collapsing magnetic field inside the motor generates a reverse voltage polarity that forward biases the diode.

The diode then provides a safe path for the inductive current to circulate temporarily, dissipating the stored magnetic energy safely and preventing high-voltage spikes from appearing across the transistor terminals.

\subsection{Power Supply Subsystem}

The robot uses two separate power supplies to isolate the logic circuitry from the motor load. A regulated 5V supply powers the sensors and transistor logic circuitry, while a separate 9V battery powers the DC motors.

Although the voltage sources are separate, both systems share a common ground connection. The common ground provides a shared voltage reference for all transistors and ensures correct base-emitter voltage operation.

\subsection{Mechanical Subsystem}

The mechanical subsystem consists of the chassis, wheels, motors, and structural support components. The robot uses a four-wheel configuration in which the front wheels are motor-driven while the rear wheels act as support wheels.

The robot operates using a differential drive mechanism. Forward motion occurs when both motors rotate simultaneously. Turning is achieved by selectively stopping one motor while the other motor continues rotating. This creates a difference in wheel velocity that changes the direction of motion.

The mechanical subsystem therefore plays a critical role in determining the stability, maneuverability, and responsiveness of the robot.

\section{Experimental Results}

The developed line-following robot was successfully tested on a floor containing a black tracking line. Experimental evaluation was performed to verify the operation of the sensing subsystem, transistor-based logic inversion stage, motor driver circuitry, and overall navigation capability of the robot.

During testing, the TCRT5000 infrared sensors successfully differentiated between the black line and the surrounding area based on infrared reflection intensity. The active-low outputs produced by the sensors were correctly inverted using 2N3904 transistor stages, allowing the motor driver subsystem to respond according to the required control logic.

The TIP122 transistor motor drivers successfully controlled the DC motors using low-current logic signals generated by the sensor circuitry. The inclusion of flyback diodes prevented instability and protected the transistor drivers from inductive voltage spikes generated during motor switching operations.


Mechanical testing revealed turning difficulties when only one motor was active. The robot initially failed to rotate effectively because the wheel friction and chassis resistance were too high for single-wheel actuation. This issue was significantly improved by increasing wheel traction using rubber bands attached around the wheels, allowing smoother differential turning behavior.

After optimization, the robot was able to move forward correctly, detect the black line reliably, and perform directional correction through differential motor control. The robot demonstrated stable line-following behavior under moderate operating speeds and consistent lighting conditions.

\section{Limitations}

Although the developed robot successfully achieved line-following functionality, several limitations were observed during implementation and testing.

The system relies entirely on discrete transistor logic and does not use a microcontroller or programmable controller. As a result, the robot lacks advanced control algorithms, adaptive speed regulation, and dynamic path prediction. This limits the precision and smoothness of line tracking at higher speeds.

The robot performance is highly dependent on surface conditions and ambient lighting. Variations in external light intensity or surface reflectivity can affect the infrared sensor readings and lead to unstable detection behavior. Proper calibration of the sensor sensitivity is therefore necessary before operation.

The use of TIP122 Darlington transistors introduces voltage drops during motor operation, reducing the effective voltage delivered to the motors. This decreases power efficiency and can reduce motor torque under load conditions.

Mechanical limitations were also observed due to the four-wheel chassis configuration. High friction between the wheels and the surface increased turning resistance, particularly during single-wheel steering operations. Although traction improvements reduced this issue, the turning mechanism remains less efficient than caster-wheel-based designs.

Furthermore, motor speed is directly dependent on battery voltage. When newer batteries were used, the robot moved too quickly and occasionally overshot turns before the sensors could respond. Since the system does not include closed-loop speed regulation, motion stability varies with battery condition.

\section{Conclusion}

This project presented the design and implementation of a transistor-based line-following robot using TCRT5000 infrared sensors, 2N3904 signal conditioning transistors, and TIP122 motor driver transistors. The system successfully demonstrated autonomous line detection and navigation without the use of microcontrollers or programmable logic devices.

The project provided practical understanding of several important electronic engineering concepts including infrared sensing, transistor switching, signal inversion, motor driver design, inductive load protection, and power supply grounding. The implementation process also highlighted the importance of proper transistor biasing, common grounding, and mechanical design considerations in embedded robotic systems.

Experimental testing confirmed that the robot could reliably detect and follow a black line on a contrasting white surface using differential motor control. Several real-world engineering challenges were encountered and resolved during development, including active-low sensor interpretation, transistor loading effects, motor switching instability, and chassis friction problems.

Despite certain limitations related to speed regulation and turning precision, the project successfully achieved its primary objective using simple and low-cost electronic components. The developed system serves as an effective educational platform for understanding the integration of analog electronics, sensing systems, and robotic motion control.

Future improvements may include the use of pulse-width modulation (PWM) speed control, microcontroller-based processing, and optimized mechanical chassis configurations to achieve smoother and more accurate navigation performance.

\begin{thebibliography}{99}

\bibitem{tcrt5000}
Vishay Semiconductors, ``TCRT5000 Reflective Optical Sensor Datasheet,'' Vishay Intertechnology, 2023.

\bibitem{tip122}
ON Semiconductor, ``TIP120/TIP121/TIP122 Darlington Transistors Datasheet,'' ON Semiconductor, 2022.

\bibitem{electronics}
R. L. Boylestad and L. Nashelsky, \textit{Electronic Devices and Circuit Theory}, 11th ed., Pearson Education, 2013.


\end{thebibliography}



\end{document}
